\chapter{Implementierung}
\label{chap:implementierung}

\section{Demo-Anwendung (Resilience Lab)}
\label{sec:demo-anwendung}
Die Webanwendung wird auf Basis von Python (FastAPI) implementiert. Sie stellt Kernendpunkte für den regulären Betrieb sowie dedizierte Schnittstellen zur Fehlerinduktion bereit:
\begin{itemize}
    \item \texttt{GET /}: Übersichtliche HTML-Statusseite mit Anzeige von Pod-Name, Anwendungsversion, Node-IP und Systemzeit.
    \item \texttt{GET /health/live}: Liveness-Prüfung zur Erkennung interner Deadlocks.
    \item \texttt{GET /health/ready}: Readiness-Prüfung zur Steuerung der Service-Teilnahme.
    \item \texttt{POST /lab/crash}: Forcierter Programmabsturz via \texttt{os.\_exit(1)}.
    \item \texttt{POST /lab/toggle-ready}: Umschaltung der Bereitschaft für kontrollierte Traffic-Aussteuerungen.
\end{itemize}

\section{Terraform-Konfiguration und k3s-Bootstrap}
\label{sec:terraform-k3s-bootstrap}
Die AWS-Infrastruktur wird über modulare Terraform-Dateien beschrieben:
\begin{itemize}
    \item \texttt{vpc.tf}: Definiert VPC, Subnetze, Internet Gateway und Routing-Tabellen.
    \item \texttt{security-groups.tf}: Konfiguriert Inbound-Regeln für SSH (Port 22), HTTP/HTTPS (Port 80/443) und die Kubernetes-API (Port 6443).
    \item \texttt{ec2.tf}: Instanziiert die VM unter Verwendung eines \texttt{cloud-init}-Skripts, welches k3s und Argo~CD installiert und das Kubeconfig-Zertifikat mit der öffentlichen IP ausstellt.
\end{itemize}

\section{Kubernetes-Manifeste mit Kustomize}
\label{sec:k8s-manifeste-kustomize}
Die Anwendungsbereitstellung wird über eine strukturierte Kustomize-Hierarchie realisiert:
\begin{itemize}
    \item \texttt{kustomize/base}: Beinhaltet Basis-Manifeste für Deployment und Service mit Standardwerten (1 Replikat, Ressourcenlimits, Health Probes).
    \item \texttt{kustomize/overlays/local}: Patcht den Service auf \texttt{NodePort} 30080 und nutzt lokale Container-Images für Entwicklungszwecke.
    \item \texttt{kustomize/overlays/aws}: Skaliert das Deployment auf 2 Replikate für Resilienztests und fügt ein Ingress-Manifest für den k3s-integrierten Traefik-Controller hinzu.
\end{itemize}

\section{Argo-CD-Anbindung und GitOps-Workflow}
\label{sec:argocd-anbindung}
Die Bereitstellung der Anwendung erfolgt über eine Argo-CD-Application-CRD. Argo~CD überwacht das Git-Repository auf dem Pfad \texttt{kustomize/overlays/aws} und synchronisiert den Zustand automatisiert (\textit{Automated Sync} \& \textit{Self-Heal}).

\section{CI-Pipeline mit GitHub Actions}
\label{sec:ci-pipeline}
Eine GitHub-Actions-Pipeline automatisiert den Integrationsprozess: Bei Änderungen am Anwendungscode werden automatisierte Tests ausgeführt, ein Container-Image gebaut und in die GitHub Container Registry (\texttt{ghcr.io}) übertragen.
